% sections/acknowledgements.tex
\cleardoublepage % Los agradecimientos comienzan en página impar
\thispagestyle{empty}

\chapter*{Agradecimientos}
\addcontentsline{toc}{chapter}{Agradecimientos}
\markboth{Agradecimientos}{Agradecimientos}

En el transcurso que he dedicado a la realización de mi tesis de maestría puedo decir con
total honestidad que no soy el mismo que comenzó con este sueño. No lo soy solo por el
tiempo que ha pasado, sino en mis formas, en mis ánimos, en mi entender acerca de lo que
quería para mí entonces y lo que quiero ahora, e incluso en la gente que me rodeó en este
importante paso. Muchos, al igual que yo, se embarcaron en este camino y, entre consejos,
incontables horas de esfuerzo, dedicación y ciertamente una cantidad gigante de obstinación
y terquedad, continúan acompañándome, siguiendo también su camino e inspirándome aun cuando
los ánimos flaqueaban. Entre ellos debo agradecer a Miguel Ángel, a Daniel Tinoco y a mi
gran amigo y hermano Jesús David Avellaneda, quien siempre estuvo preguntando y alentándome
a continuar. Hay otras personas que empezaron esto a mi lado; su presencia, su alegría y su
esfuerzo fueron parte de mi comienzo aquí. Quiero agradecer a Carolina Carrillo Rodríguez
por ser parte de esa inspiración.

No tengo más que palabras de agradecimiento hacia mi tutor, el doctor Herbert Vinck Posada,
sin cuya infinita paciencia y apoyo este trabajo no habría sido posible. Supo esperar
cuando el tiempo no me alcanzaba y animarme cuando el trabajo parecía no avanzar, sin dejar
de creer que llegaría hasta aquí. Por su generosidad, su guía y su confianza, le guardo una
gratitud que estas líneas apenas alcanzan a expresar.

Desde que era niño y decidía quedarme en la biblioteca leyendo, buscando algo sin saber qué,
en lugar de salir con los otros niños, una presencia estuvo siempre guiándome e impulsando
mi curiosidad. Me llevó de la mano incluso sin saber cuál era el camino que debíamos seguir,
sin darse por vencida por ello, y me ayudó con las tareas más complicadas hasta el día en que
los conocimientos la sobrepasaron, pero no las ganas ni el querer verme triunfar. Por su
empeño en cumplir una palabra que jamás olvidaré, «así tenga que lanzarlo del pelo por la
pared, ha de ir a la universidad», la persona más importante a la que debo agradecer es a mi
madre, cuyo esfuerzo infinito me ha traído a este punto de mi vida. ¡Esta es tu victoria!
